\begin{mistake}{2026-05-02}{05}

\begin{problemcontent}
辨析：含有第一类间断点（如跳跃间断点、可去间断点）的函数 \(f(x)\) 在包含该间断点的区间内，是否具有原函数？是否必定不可积？
\end{problemcontent}

\begin{solutioncontent}
\textbf{1. 是否有原函数：必定没有。}
原函数的要求是存在 \(F(x)\) 使得处处有 \(F'(x) = f(x)\)。
根据\textbf{达布定理}（导数介值定理），任何导函数都不会出现第一类间断点。既然 \(f(x)\) 存在第一类间断点，它就绝对不可能成为某个函数的导函数，故必无原函数。

\textbf{2. 是否可积：必然可积。}
定积分的本质是面积极限。根据\textbf{黎曼可积的充分条件}，只要函数在闭区间上\textbf{有界}且只有\textbf{有限个间断点}，定积分就必然存在。第一类间断点只是图形断开，未趋于无穷（满足有界性），因此完全可积。

\textbf{结论}：其变限积分 \(\Phi(x) = \int_a^x f(t)\,\mathrm{d}t\) 存在且连续，但在间断点处不可导，所以 \(\Phi(x)\) 也不是它的原函数。
\end{solutioncontent}

\mistakeknowledge{
\begin{itemize}
  \item \textbf{核心定理}：原函数存在定理（连续必有原函数）；可积充分条件（有界且间断点有限必可积）。
  \item \textbf{易错点}：将"变限积分函数"等同于"原函数"，混淆"不可积"与"无原函数"。带有第一类间断点的函数是"可积但无原函数"的典型代表。
  \item \textbf{关联知识}：既然无整体原函数，求定积分时不可直接套用牛顿-莱布尼茨公式，必须在间断点处拆分区间分段计算。
\end{itemize}}

\end{mistake}
